\section{Module Selection}
\label{sec:proposed:modules}

The brainstorming phase produced a catalogue of nineteen
analytical modules grouped under five lenses: Network Overview,
Route Explorer, Stop Analysis, Financial \& Operations, and
Metro Connectivity. Each module is grounded in one of the
literature sources reviewed in Chapter~\ref{chapter:related}: the
QoS framework of Srisakaokul~\textit{et
al.}~\cite{srisakaokul2023}, the FuzzyCut, GroupSet, and Gridify
primitives of Liu~\cite{liu2022}, and the analytical
specification of the advisor~\cite{thoai2025}. Not every module
is supported by the data available locally. This section screens
each module against the schema (Section~\ref{sec:proposed:schema})
and the EDA findings (Section~\ref{sec:proposed:eda}), and assigns
each one to one of three categories.

\begin{description}
\item[$\checkmark$ Adopted.] The module is supported by the local
    sample at full coverage and is implemented in the dashboard
    as specified.
\item[$\sim$ Adopted with sample disclaimer.] The module is
    supported but the local data limits it to nine GPS routes
    and a single day. The dashboard surfaces the limitation in the
    UI and the analyst is invited to interpret results
    accordingly.
\item[$\times$ Deferred.] The module requires data not present in
    the local release (e.g., the metro timetable, multi-day GPS
    coverage). It is documented as planned future work in
    Chapter~\ref{chapter:conclusion}.
\end{description}

The full screening result is given in
Table~\ref{tab:proposed:modules}; the subsections that follow
walk through the rationale for each lens.

\begin{table}[H]
\centering
\caption{Analytics-module screening against the local data sample.}
\label{tab:proposed:modules}
\renewcommand{\arraystretch}{1.20}
\begin{tabularx}{\linewidth}{p{3.5cm}p{2cm}X}
\hline
\textbf{Module} & \textbf{Status} & \textbf{Reason} \\
\hline
\multicolumn{3}{l}{\emph{Network Overview}} \\
KPI strip                & $\checkmark$ & AFC supplies boardings, revenue, ghost\%, busiest route. Network-wide QoS-3 deferred (needs schedule + GPS at scale). \\
Demand cluster (GroupSet) & $\checkmark$ & Hourly demand per route across 90 days is sufficient for GroupSet. \\
Route health treemap     & $\checkmark$ & Ghost-rate substituted for QoS-3 as the colour scale. \\
Operator benchmarking    & $\checkmark$ & 10 companies $\times$ 128 routes provides enough density for the scatter. \\
\hline
\multicolumn{3}{l}{\emph{Route Explorer}} \\
QoS-1 reliability        & $\sim$       & Schedule data exists for 9 GPS routes; trend visible only on those. \\
QoS-2 accessibility      & $\sim$       & Aggregate on-path rate of 25.1\% is low because of urban-canyon noise; relative interpretation only. \\
QoS-3 availability       & $\sim$       & Headway deviation requires schedule; 9 GPS routes only. \\
Speed profile (FuzzyCut) & $\sim$       & GPS speed exists on 9 routes. Noise filter applied at HMM stage. \\
Trip Gantt timeline      & $\sim$       & 9 routes, 1 day. \\
Hourly demand            & $\checkmark$ & AFC-driven; works for any of 128 routes. \\
Holiday vs.\ normal      & $\sim$       & Adapted: AFC-side comparison around Tet (2026-02-17) replaces the GPS-side parade-period comparison until full data lands. \\
\hline
\multicolumn{3}{l}{\emph{Stop Analysis}} \\
Stop detail panel        & $\checkmark$ & Boardings + routes from AFC; dwell as enrichment when stop is in GPS sample. \\
Spatial accessibility (Gridify) & $\checkmark$ & Stops master with lat/lng covers all 4{,}710 HCMC stops. \\
Multi-stop comparison    & $\sim$       & Punctuality dimension is GPS-only; other dimensions (boardings, dwell, routes served) work everywhere. \\
\hline
\multicolumn{3}{l}{\emph{Financial \& Operations}} \\
Ghost trip analysis      & $\checkmark$ & 17.8\% near-ghost rate at network scale; per-route ranking with $>40\%$ outliers. \\
Revenue \& ticket breakdown & $\checkmark$ & AFC + payments fully populated. \\
Revenue efficiency (FuzzyCut) & $\checkmark$ & Revenue-per-trip distribution exists for all 128 routes. \\
\hline
\multicolumn{3}{l}{\emph{Metro Connectivity}} \\
Feeder catchment rings    & $\checkmark$ & Geometric: stops master + metro station coordinates. \\
Transfer timing analysis  & $\times$    & Requires the metro Line 1 timetable, which is not in the data release. \\
QoS-3 at metro interfaces & $\times$    & Requires schedule + GPS coverage of metro-adjacent stops. Both deferred. \\
Feeder headway alignment (GroupSet) & $\times$ & Same dependency on schedule data. \\
\hline
\end{tabularx}
\end{table}

\subsection{Network Overview}

All four Network modules are adopted. The KPI strip drops the
network-wide QoS-3 indicator that the original brainstorm proposed
because the schedule-side data exists only for nine routes and
a single day; with that exclusion, the strip is fully populated
from AFC-side aggregates alone (boardings, revenue, ghost rate,
busiest route, and the demographic split). The treemap is sized
by total boardings, but it is coloured by ghost rate rather than by
QoS-3 because the former is computable for all 128 routes.
The GroupSet demand-pattern clusterer operates on AFC hourly
demand per route across the 90-day window, which is the data
density the technique was originally designed for.

\subsection{Route Explorer}

All three QoS indicators are adopted with a sample disclaimer.
QoS-2 is interpreted relatively rather than absolutely because of
the 25.1\% baseline on-path rate. The Speed Profile FuzzyCut
corridor and the Trip Gantt timeline are likewise restricted to
the nine GPS routes. Hourly demand from AFC is unaffected by the
GPS limitation and is available for any of the 128 routes.
The Holiday vs.\ Normal comparison is adapted: instead of the
GPS-side Apr 30 parade period (which falls outside the local
sample), the lens compares AFC demand around Tet
(2026-02-17, the Lunar New Year) to a normal weekday baseline.
The original GPS-side comparison is queued for the full-dataset
release.

\subsection{Stop Analysis}

The Stop Detail panel and the Gridify accessibility view are
adopted at full coverage. The Multi-Stop Comparison radar is
adopted with the punctuality dimension surfaced only when both
selected stops fall within the nine-route GPS sample;
boardings, dwell, and routes-served are computable for every
stop in the master.

\subsection{Financial and Operations}

All three modules are adopted. The Ghost Trip analysis is
particularly well-supported: the 17.8\% network-wide near-ghost
rate provides a defendable headline number, and per-route ranking
identifies several routes above 40\%. The Revenue Efficiency
FuzzyCut applies the same primitive as the Route Explorer's Speed
Profile, which is one of the design wins of using literature
primitives: the analyst learns the FuzzyCut interaction once and
reuses it for two unrelated quantities.

\subsection{Metro Connectivity}

The Metro lens is the most data-dependent of the five.
The Feeder Catchment rings are computable from the stops master and
the published metro station coordinates. The remaining three modules
(Transfer Timing, QoS-3 at Metro Interfaces, Feeder Headway
Alignment) all require either the metro Line 1 timetable, multi-day
GPS coverage of metro-adjacent stops, or both. They are documented
as planned next steps in Chapter~\ref{chapter:conclusion} and the
lens icon carries an ``in progress'' badge in the dashboard until
the dependencies are resolved.
